% ---------------------------------------------------------------------------
% content/conclusion.tex -- an unnumbered chapter (see main.tex for how it is
% added to the table of contents and anchored for the Hebrew TOC).
% ---------------------------------------------------------------------------

This template takes care of the mechanics that make a BGU thesis awkward to typeset: an English body and a Hebrew back part in two documents, compiled by two engines, joined into a single PDF whose tables of contents, page numbers and links agree.
The conventions are few.
Chapters are listed in \texttt{main.tex} with an anchor and a label, sections are written with \verb|\linkedsection|, and any theorem, figure or definition the Hebrew text refers to gets an anchor next to its label.
On the Hebrew side, \verb|\en| marks Latin fragments, \verb|\hcite| cites, and \verb|\enref| links to an English result and prints its number.
The script compiles both parts until they agree and merges them, and its merge step tells you if a link has nowhere to go.

To turn the template into your thesis, fill in \texttt{thesis\_info.tex}, replace the chapters in \texttt{content/} and the abstract and acknowledgements in \texttt{Main\_pages/}, write the Hebrew summary and abstract in \texttt{hebrew\_pages.tex}, and keep its table of contents in step with your chapters.
Include the affidavit if your department requires it, and decide whether the thesis will be printed double-sided or read on screen, which determines whether chapters should open on odd pages.
Before submission, build from a clean \texttt{build/} folder, make sure the merge step prints no warning, and open the PDF in the viewer your readers will use to click a few links in both directions.
